\chapter{Manual de usuario}
\label{anx:manual}

Este anexo ofrece una guía paso a paso de las operaciones principales
de \dataqual\ desde la perspectiva del usuario final.  Se asume que el
sistema está desplegado y accesible en \texttt{http://localhost:3000}.
La interfaz está disponible en español e inglés y el idioma se puede
alternar en cualquier momento desde el selector \texttt{ES~/~EN}
situado en la parte inferior de la barra lateral.

\section*{Requisitos previos}

\begin{enumerate}
  \item Docker Desktop instalado y en ejecución.
  \item Repositorio clonado: \texttt{git clone <url-repositorio>}.
  \item Fichero \texttt{backend/.env} configurado con las variables de
    entorno, a partir de la plantilla
    \texttt{backend/.env.example}.
\end{enumerate}

\section*{Arranque del sistema}

\begin{verbatim}
cd TFG_CALIDAD_DATOS_IA
docker-compose up --build   # Primera vez (compila imágenes)
docker-compose up           # Arranques posteriores
\end{verbatim}

El sistema tarda aproximadamente 30--60 segundos en estar completamente
operativo.  Una vez listos los servicios, acceder a
\texttt{http://localhost:3000}.

\section*{Navegación general}

La interfaz se organiza en torno a tres elementos: una barra lateral
persistente que agrupa la navegación principal, una zona central de
contenido y un encabezado contextual con \emph{breadcrumbs} que
indica la posición actual (proyecto, \emph{dataset} o evaluación
activos).

La barra lateral se divide en dos secciones.  La sección principal
contiene los apartados \textbf{Dashboard}, \textbf{Proyectos},
\textbf{Datasets} y \textbf{Análisis}; los tres últimos se despliegan
en accesos directos a la vista de listado, a la acción de creación o
carga, y a las entradas recientes (los proyectos del usuario aparecen
como atajos directos bajo \textbf{Proyectos}).  La sección de
\textbf{Configuración} agrupa el perfil del usuario y la gestión de
plantillas de métricas reutilizables.  En la parte inferior se sitúan
el selector de idioma \texttt{ES~/~EN}, el indicador del usuario
activo y el botón de cierre de sesión.  La barra lateral se puede
colapsar a un panel de iconos pulsando el \emph{chevron} de su
cabecera, lo que maximiza el área útil para el contenido.

\section*{Cuadro de mando}

Tras iniciar sesión el sistema redirige al cuadro de mando
(\texttt{/dashboard}), que ofrece una visión global del estado de
calidad de todos los proyectos del usuario.  La pantalla agrupa
varias secciones de información:

\begin{itemize}
  \item \textbf{Indicadores clave (KPI):} número de proyectos
    activos, \emph{datasets} en uso y puntuación media de los
    últimos análisis completados.
  \item \textbf{Distribución de \emph{Quality Gates}:} desglose
    visual del estado del último análisis de cada proyecto, con
    contadores de \textbf{Aprobado}, \textbf{Advertencia} y
    \textbf{Fallido}.
  \item \textbf{Requiere atención:} listado de los proyectos cuyo
    \emph{Quality Gate} ha fallado, está en advertencia, o cuya
    puntuación de calidad se encuentra por debajo del 70\,\%.
  \item \textbf{Historial de salud:} gráfico de barras con el
    histórico de análisis recientes por proyecto, accesible
    mediante \emph{scroll} horizontal para revisar la evolución a
    lo largo del tiempo.
\end{itemize}

Si el usuario aún no ha creado ningún proyecto, el cuadro de mando
muestra un estado vacío con una llamada a la acción \textbf{Crear
proyecto}.

\section*{Flujo principal de uso}

\subsection*{1. Registro e inicio de sesión}

La autenticación se realiza desde la pantalla unificada
\texttt{/auth}, que ofrece dos pestañas: \textbf{Iniciar sesión}
(activa por defecto) y \textbf{Registrarse}.

\begin{enumerate}
  \item Para crear una cuenta nueva, seleccionar la pestaña
    \textbf{Registrarse} (o acceder directamente a
    \texttt{/auth?mode=register}) e introducir nombre de usuario,
    correo electrónico y contraseña.  Tras el registro, el sistema
    inicia sesión automáticamente y redirige al cuadro de mando.
  \item Para los inicios de sesión sucesivos, introducir el usuario
    o correo y la contraseña en la pestaña \textbf{Iniciar sesión}.
    El \emph{token} \ac{JWT} se almacena en \texttt{localStorage} y
    se inyecta automáticamente en todas las peticiones a la \ac{API}
    mediante el cliente Axios; la renovación del \emph{token} de
    acceso es transparente para el usuario gracias al \emph{refresh
    token} de larga duración.
\end{enumerate}

\begin{figure}[htbp]
  \centering
  
\includegraphics[width=0.85\textwidth]{img/manual-auth.png}
  \caption{Pantalla unificada de autenticación con pestañas de
    inicio de sesión y registro}
  \label{fig:manual-auth}
\end{figure}

\subsection*{2. Crear un proyecto}

Los proyectos son la unidad organizativa de \dataqual: agrupan los
\emph{datasets} cargados, las métricas configuradas, las
evaluaciones ejecutadas y el \emph{Quality Gate} del conjunto.

\begin{enumerate}
  \item Pulsar \textbf{Nuevo proyecto} desde el cuadro de mando,
    desde el listado de proyectos (\texttt{/projects}) o desde el
    acceso directo de la barra lateral.  La aplicación abre un
    asistente de tres pasos:
    \begin{enumerate}
      \item \textbf{Información básica:} nombre del proyecto y
        descripción opcional.
      \item \textbf{Plantilla y métricas:} elección entre aplicar
        una plantilla del sistema (que preconfigura un conjunto
        coherente de métricas y umbrales) o una configuración
        personalizada activando individualmente las métricas
        deseadas.  Los parámetros detallados se pueden refinar
        posteriormente desde la página del proyecto.
      \item \textbf{Resumen y confirmación:} revisión final del
        proyecto antes de su creación.
    \end{enumerate}
  \item Al confirmar, el proyecto aparece en el listado con estado
    \emph{Sin evaluaciones} y queda disponible como atajo en la
    barra lateral para acceso rápido en sesiones posteriores.
\end{enumerate}

\begin{figure}[htbp]
  \centering
  \includegraphics[width=0.85\textwidth]{img/manual-new-project.png}
  \caption{Asistente de creación de un nuevo proyecto en tres pasos}
  \label{fig:manual-new-project}
\end{figure}

\subsection*{3. Cargar un \emph{dataset}}

La carga de \emph{datasets} sigue un asistente de tres pasos
accesible desde \texttt{/datasets/upload}, desde el botón
\textbf{Añadir dataset} en la pestaña \emph{Datasets} de un proyecto,
o desde el acceso directo \textbf{Subir dataset} en la barra
lateral.

\begin{enumerate}
  \item \textbf{Selecciona un proyecto:} elegir el proyecto al que
    se asociará el \emph{dataset}.  Si la carga se inicia desde la
    página de un proyecto concreto, este paso se completa de forma
    automática.
  \item \textbf{Añade un archivo:} arrastrar y soltar un fichero
    CSV en la zona de carga o pulsar para abrir el selector de
    archivos.  El sistema solo admite ficheros con extensión
    \texttt{.csv} y un tamaño máximo de 10\,MB, y aplica una
    estrategia de \emph{fallback} para codificaciones y
    delimitadores poco habituales descrita en la
    sección~\ref{sec:impl-ingesta}.
  \item \textbf{Revisa y confirma:} introducir el nombre del
    \emph{dataset}, una descripción opcional y, opcionalmente, una
    etiqueta de versión (por ejemplo \texttt{v2.0},
    \texttt{producción} o \texttt{baseline}).  En la misma pantalla,
    el panel \emph{Privacidad de datos} permite marcar las columnas
    sensibles (\ac{PII}) cuyos valores se enmascararán
    automáticamente en todas las salidas posteriores del sistema
    (logs, \emph{issues}, vistas previas, perfilado).
\end{enumerate}

Al confirmar la carga, el sistema valida el fichero, extrae el
esquema y las estadísticas descriptivas básicas, almacena el binario
en MinIO y registra los metadatos en PostgreSQL.

\begin{figure}[htbp]
  \centering
  \includegraphics[width=0.85\textwidth]{img/manual-upload.png}
  \caption{Paso 2 del asistente de carga: zona de \emph{drag-and-drop}
    con nombre y descripción del \emph{dataset}}
  \label{fig:manual-upload}
\end{figure}

\begin{figure}[htbp]
  \centering
  \includegraphics[width=0.85\textwidth]{img/enmascarado-pii.png}
  \caption{Paso 3 del asistente de carga: resumen del \emph{dataset} y
    selección de columnas \ac{PII} para enmascaramiento automático}
  \label{fig:manual-pii}
\end{figure}

\subsection*{4. Configurar métricas del proyecto}

Cada proyecto mantiene su propia configuración de métricas, accesible
desde la pestaña \textbf{Métricas} de la página del proyecto o desde
el botón \textbf{Configurar métricas}, que abre la página dedicada
\texttt{/metrics/configure/[id]}.  El configurador presenta tres
pestañas:

\begin{enumerate}
  \item \textbf{Métricas disponibles:} catálogo de las siete
    métricas evaluables del sistema (completitud, unicidad,
    exactitud sintáctica, consistencia lógica, equilibrio de clases,
    actualidad y diversidad).  Cada métrica incluye una breve
    descripción de lo que mide y un botón para activarla en el
    proyecto.
  \item \textbf{Métricas seleccionadas:} para cada métrica activa,
    se pueden ajustar sus parámetros específicos (umbrales, pesos,
    columnas objetivo, patrones de validación para la exactitud
    sintáctica, reglas \texttt{IF--THEN} para la consistencia
    lógica, etc.).
  \item \textbf{Plantillas:} aplicar una plantilla predefinida del
    sistema o una creada por el usuario (véase paso 10) para
    reutilizar configuraciones entre proyectos similares sin tener
    que repetir el ajuste manual.
\end{enumerate}

Pulsar \textbf{Guardar configuración} aplica los cambios y devuelve
al usuario a la página del proyecto.

\begin{figure}[htbp]
  \centering
  \includegraphics[width=0.85\textwidth]{img/manual-metrics.png}
  \caption{Configurador de métricas (pestaña \emph{Métricas seleccionadas}):
    tarjetas de las métricas activas con acceso a configuración detallada
    y opción de guardar como plantilla reutilizable}
  \label{fig:manual-metrics}
\end{figure}

\subsection*{5. Ejecutar una evaluación}

Una vez configuradas las métricas y cargado al menos un
\emph{dataset}, el usuario puede lanzar un análisis desde la página
del proyecto pulsando el botón \textbf{Ejecutar análisis}.

\begin{enumerate}
  \item Seleccionar el \emph{dataset} (y opcionalmente la versión
    concreta) que se desea evaluar.
  \item El sistema encola la tarea en Celery y muestra una barra
    de progreso en tiempo real (0--100\,\%) actualizada mediante
    \emph{polling} cada pocos segundos.  Si Celery no estuviera
    disponible, el servicio híbrido (\texttt{HybridEvaluationService})
    ejecuta el análisis de forma síncrona como mecanismo de
    \emph{fallback}, garantizando que la evaluación se complete.
  \item Al finalizar, el sistema registra los resultados como un
    \texttt{AnalysisRun} inmutable y muestra el \emph{Quality
    Score} obtenido junto con el veredicto del \emph{Quality Gate}.
    Si la evaluación se quedara atascada por un tiempo prolongado,
    el \emph{watchdog} de evaluaciones la detectaría y marcaría
    como fallida automáticamente.
\end{enumerate}

\begin{figure}[htbp]
  \centering
  \includegraphics[width=0.85\textwidth]{img/manual-evaluation.png}
  \caption{Pestaña \emph{Evaluaciones} de un \emph{dataset}: evaluación
    recién lanzada con estado \emph{Queued} y botón \emph{Run evaluation}
    en curso, junto a una evaluación anterior con resultado PASSED}
  \label{fig:manual-evaluation}
\end{figure}

\subsection*{6. Consultar resultados}

Los resultados de los análisis son accesibles desde tres puntos
de entrada: la pestaña \textbf{Análisis} de la página del proyecto
(histórico de análisis del proyecto), el listado global en
\texttt{/evaluations} (histórico cruzado de todos los proyectos),
y la pestaña \textbf{Evaluaciones} de la página del \emph{dataset}.

\begin{enumerate}
  \item Pulsar sobre un análisis completado para abrir su detalle.
  \item El indicador circular (\emph{gauge}) muestra el
    \emph{Quality Score} global (0--100), acompañado del veredicto
    del \emph{Quality Gate} (\texttt{PASSED}, \texttt{WARNING} o
    \texttt{FAILED}) y de un desglose por dimensiones de calidad.
  \item Las pestañas dedicadas a cada métrica detallan el resultado
    por columna y permiten profundizar en los \emph{issues}
    detectados, indicando severidad (crítica, alta, media, baja),
    columna afectada, descripción del problema y, cuando procede,
    el porcentaje de filas afectadas.
  \item Los \emph{badges} que acompañan a cada \emph{issue} indican
    su estado respecto a la evaluación anterior:
    \textbf{Nuevo}, \textbf{Recurrente} o, en la pestaña de cambios,
    \textbf{Corregido}.  Esta trazabilidad se sustenta en el sistema
    de \emph{fingerprinting} descrito en la
    sección~\ref{sec:impl-fingerprinting}, que asigna un \emph{hash}
    determinista a cada problema.
\end{enumerate}

\begin{figure}[htbp]
  \centering
  \includegraphics[width=0.85\textwidth,height=0.85\textheight,keepaspectratio]{img/manual-results.png}
  \caption{Detalle de un análisis: \emph{Quality Score}, métricas
    e \emph{issues} con \emph{badges} de trazabilidad
    entre ejecuciones}
  \label{fig:manual-results}
\end{figure}

\subsection*{7. Explorar datos (\acs{EDA})}

El módulo de perfilado interactivo es accesible desde la pestaña
\textbf{Perfilado} de cualquier \emph{dataset}
(\texttt{/datasets/[id]}).  Concentra en una única pantalla las
visualizaciones habituales del análisis exploratorio de datos:

\begin{enumerate}
  \item Histogramas y diagramas de caja (\emph{boxplots}) por
    columna, con vistas separadas para columnas numéricas y
    categóricas.
  \item Matriz de correlación, alternable entre los coeficientes
    de Pearson (relaciones lineales entre variables numéricas) y
    Spearman (relaciones monótonas, incluidas variables
    ordinales).
  \item Detección de valores atípicos (\emph{outliers}) con factor
    \ac{IQR} ajustable interactivamente para explorar distintos
    niveles de tolerancia.
  \item Análisis de duplicados a nivel de fila completa y de
    variabilidad por columna, útil para identificar columnas
    quasi-constantes o claves duplicadas.
\end{enumerate}

Las columnas previamente marcadas como sensibles aparecen con sus
valores ofuscados en todas las visualizaciones, garantizando que la
exploración exploratoria no exponga información personal.

\begin{figure}[htbp]
  \centering
  \includegraphics[width=0.85\textwidth,height=0.85\textheight,keepaspectratio]{img/manual-eda.png}
  \caption{Módulo de perfilado interactivo (\ac{EDA}): histogramas,
    diagramas de caja, correlaciones y detección de
    \emph{outliers}}
  \label{fig:manual-eda}
\end{figure}

\subsection*{8. Seguir la evolución temporal del proyecto}

La pestaña \textbf{Analysis history} de la página de un proyecto
concentra el historial completo de análisis ejecutados, con una
perspectiva de seguimiento longitudinal:

\begin{itemize}
  \item \textbf{Gráfico de evolución:} serie temporal de la puntuación
    de calidad por \emph{dataset}, con indicadores de tendencia positiva
    o negativa entre ejecuciones consecutivas.
  \item \textbf{Filtros por dataset:} posibilidad de aislar la evolución
    de un \emph{dataset} concreto o verlos todos superpuestos.
  \item \textbf{Tabla de ejecuciones:} lista paginada de todos los
    \texttt{AnalysisRun} del proyecto con fecha, \emph{dataset},
    puntuación, número de \emph{issues} y estado del \emph{Quality Gate}.
\end{itemize}

\begin{figure}[htbp]
  \centering
  \includegraphics[width=0.85\textwidth]{img/analysis-history.png}
  \caption{Pestaña \emph{Analysis history} de un proyecto: gráfico de
    evolución temporal de la puntuación de calidad por \emph{dataset}
    y tabla de ejecuciones con estado de \emph{Quality Gate}}
  \label{fig:manual-analysis-history}
\end{figure}

\subsection*{9. Configurar un \emph{Quality Gate}}

Cada proyecto puede definir su propio \emph{Quality Gate} desde la
pestaña \textbf{Quality Gate} de la página del proyecto.  El editor
ofrece dos modos de configuración complementarios:

\begin{enumerate}
  \item \textbf{Preajustes de rigor:} botones de selección rápida
    para los perfiles \textbf{Estricto}, \textbf{Moderado} y
    \textbf{Permisivo}, que aplican un conjunto coherente de
    umbrales sin necesidad de ajustarlos individualmente.  Útiles
    como punto de partida.
  \item \textbf{Configuración manual:} ajuste fino de los tres
    umbrales que determinan el veredicto del \emph{gate}:
    \begin{itemize}
      \item \textbf{Score mínimo} (defecto 70): puntuación de
        calidad global por debajo de la cual el \emph{Quality Gate}
        falla.
      \item \textbf{Máx.\ \emph{issues} críticos} (defecto 0):
        número máximo de \emph{issues} de severidad crítica que se
        permiten antes de fallar el \emph{gate}.
      \item \textbf{Máx.\ \emph{issues} nuevos} (defecto 10):
        número máximo de \emph{issues} nuevos respecto a la línea
        base permitidos.
    \end{itemize}
\end{enumerate}

El veredicto del \emph{Quality Gate} adopta tres estados, mostrados
tanto en la página del análisis como en el cuadro de mando:
\texttt{PASSED} (verde), \texttt{WARNING} (amarillo, dentro del
margen de advertencia configurado) y \texttt{FAILED} (rojo).  Las
opciones \textbf{Guardar cambios} y \textbf{Restaurar valores por
defecto} controlan la persistencia de la configuración.

\begin{figure}[htbp]
  \centering
  \includegraphics[width=0.85\textwidth]{img/manual-quality-gate.png}
  \caption{Pestaña \emph{Quality Gate} de un proyecto: condición de
    aprobación activa, umbrales configurables con deslizadores y
    etiquetas de nivel de rigor por umbral}
  \label{fig:manual-qg}
\end{figure}

\subsection*{10. Gestionar versiones de un \emph{dataset}}

La pestaña \textbf{Versiones} de cualquier \emph{dataset} muestra el
árbol completo de versiones del fichero, con las relaciones
padre-hijo entre cargas sucesivas y el indicador
\texttt{is\_latest} sobre la versión activa (la representación
visual se ilustra en la Figura~\ref{fig:dataset-lineage} del
Capítulo~\ref{chap:desarrollo}).  Desde esta vista se puede:

\begin{itemize}
  \item Subir una nueva versión del \emph{dataset}: la nueva carga
    se vincula automáticamente como hija de la versión seleccionada
    y se convierte en la nueva \texttt{is\_latest}.
  \item Asignar una etiqueta descriptiva a una versión
    (\texttt{v2.0}, \texttt{producción}, \texttt{baseline}\ldots)
    para facilitar su identificación posterior.
  \item Comparar dos versiones del mismo \emph{dataset} desde
    \texttt{/datasets/compare} para ver diferencias en esquema,
    estadísticas descriptivas y resultados de evaluación.
\end{itemize}

\subsection*{11. Plantillas de métricas reutilizables}

La sección \texttt{/settings/templates}, accesible desde el
apartado \textbf{Configuración} de la barra lateral, permite
gestionar plantillas del usuario que encapsulan una configuración
completa de métricas (selección, parámetros y umbrales).  Una vez
guardada, una plantilla puede aplicarse durante el paso
\emph{Plantilla y métricas} del asistente de creación de un proyecto
nuevo, o desde la pestaña \emph{Plantillas} del configurador de
métricas de un proyecto existente.  Esta funcionalidad ahorra
configuración manual repetitiva entre proyectos similares y
favorece la consistencia de criterios de calidad entre equipos.

\section*{Monitorización del sistema}

Más allá de la interfaz de usuario, dos paneles auxiliares permiten
supervisar el estado del \emph{stack} desde el navegador:

\begin{itemize}
  \item \textbf{Flower} (monitor de Celery):
    \texttt{http://localhost:5555}.  Visualiza en tiempo real los
    \emph{workers} activos, las tareas en curso y completadas, sus
    tiempos de ejecución y el histórico de fallos, y permite
    reintentar manualmente tareas concretas.
  \item \textbf{MinIO Console}: \texttt{http://localhost:9001}
    (usuario y contraseña configurados en \texttt{.env}).  Permite
    inspeccionar los \emph{buckets} de almacenamiento y verificar
    que los \emph{datasets} se han subido correctamente.
\end{itemize}

\section*{Apagado del sistema}

\begin{verbatim}
docker-compose down           # Apagar (conserva datos en volúmenes)
docker-compose down -v        # Apagar y eliminar volúmenes (borra datos)
\end{verbatim}


% Local Variables:
%  coding: utf-8
%  mode: latex
%  mode: flyspell
%  ispell-local-dictionary: "castellano8"
% End:
